\section{Worked Examples: 评估如何真的算出来}

这一讲最容易停留在概念层。为了避免这种情况，这里把几个最常见的评估任务拆成具体算例。

\subsection{例 1：多任务知识评估怎么汇总}

假设你有三个子任务，每个任务都是 10 道题，模型分别答对 8、6、9 题。最直接的 micro average 是
\[
\frac{8+6+9}{30} = 0.77.
\]
如果改成按任务平均的 macro average，则是
\[
\frac{0.8+0.6+0.9}{3} = 0.766\overline{6}.
\]
这两个数接近，但当每个子任务样本数不同的时候，它们会明显不一样。

\begin{table}[H]
\centering
\begin{tabular}{p{2.8cm}p{2.6cm}p{2.6cm}p{2.6cm}}
\toprule
\textbf{子任务} & \textbf{题数} & \textbf{答对} & \textbf{准确率} \\
\midrule
task A & 10 & 8 & 80\% \\
task B & 10 & 6 & 60\% \\
task C & 10 & 9 & 90\% \\
\bottomrule
\end{tabular}
\caption{一个简单的多任务汇总例子}
\end{table}

\begin{knowledgebox}{macro 和 micro 的差异}
micro average 让大任务占更大权重，macro average 让每个任务权重相同。一个偏向“总体通过率”，一个偏向“任务均衡性”。
\end{knowledgebox}

\subsection{例 2：pairwise 胜率怎么理解}

假设某个模型在 20 场对决里赢了 13 场，输 7 场，那么原始 win rate 是 65\%。
但这不是完整信息，因为：
\begin{itemize}
    \item 对手是谁很重要。
    \item 哪些 prompt 被抽中很重要。
    \item 置信区间有多宽也很重要。
\end{itemize}

如果你把 pairwise 结果想成二项分布，那么一个粗略的标准误差可以写成
\[
\sqrt{\frac{p(1-p)}{n}},
\]
其中 $p=0.65, n=20$。这说明样本还不算大，波动可能不小。

\begin{importantbox}{ELO 不是神谕}
ELO 的价值是把一堆 pairwise 关系压缩成一个可比较排序，但它仍然依赖采样分布、比较对象和更新规则。排行榜看起来是一个数，实际是一整套实验设计。
\end{importantbox}

\subsection{例 3：安全评估里的混淆矩阵}

安全任务经常可以被写成一个二分类问题：模型的响应要么安全，要么不安全。此时就会有四种情况：
\begin{itemize}
    \item 真安全且被判安全。
    \item 真安全但被误判不安全。
    \item 真不安全但被放过。
    \item 真不安全且被正确拦截。
\end{itemize}

\begin{figure}[H]
\centering
\begin{tikzpicture}[node distance=0.8cm, every node/.style={font=\small, align=center}]
\node[draw, rounded corners, fill=green!10, minimum width=2.8cm, minimum height=0.8cm] (a) {true safe};
\node[draw, rounded corners, fill=red!10, right=1.1cm of a, minimum width=2.8cm, minimum height=0.8cm] (b) {true unsafe};
\node[draw, rounded corners, fill=green!20, below=0.9cm of a, minimum width=2.8cm, minimum height=0.8cm] (c) {pred safe};
\node[draw, rounded corners, fill=red!20, below=0.9cm of b, minimum width=2.8cm, minimum height=0.8cm] (d) {pred unsafe};
\draw[-{Latex[length=3mm]}, thick] (a) -- (c);
\draw[-{Latex[length=3mm]}, thick] (a) -- (d);
\draw[-{Latex[length=3mm]}, thick] (b) -- (c);
\draw[-{Latex[length=3mm]}, thick] (b) -- (d);
\end{tikzpicture}
\caption{安全评估可以被组织成一个混淆矩阵问题，但真正的难点在于标签定义和上下文依赖}
\end{figure}

\begin{warningbox}{安全类指标最容易出错的地方}
把“拒答率高”误认为“安全性高”。如果模型过度拒答，它可能更安全，也可能只是更没用；如果模型过度配合，它可能更好用，也可能更危险。必须看任务定义。
\end{warningbox}

\subsection{例 4：模型研究和产品决策看的是不同东西}

如果你是研究者，往往关心一个改动是否真正提升了能力边界；如果你是产品经理，你可能更关心用户留存、价格和响应速度。

\begin{table}[H]
\centering
\begin{tabular}{p{3.0cm}p{4.3cm}p{5.4cm}}
\toprule
\textbf{角色} & \textbf{主要目标} & \textbf{常用评估} \\
\midrule
researcher & 可发表、可复现、可外推的改进 & benchmark, ablation, scaling law \\
product builder & 更好的实际体验与成本控制 & arena, traffic, latency, cost \\
policy/safety team & 风险可控、可审计 & curated safety suites, red teaming \\
\bottomrule
\end{tabular}
\caption{同一模型在不同组织里的评估目标并不一样}
\end{table}

\begin{knowledgebox}{为什么这一点重要}
如果评估目标混淆，团队会把资源投到错误的优化方向上。研究上的胜利不一定是产品上的胜利，产品上的胜利也不一定是安全上的胜利。
\end{knowledgebox}

